% !TEX root = main.tex

\section{Metodologia}\label{sec:proposed_solution}

---- IDEIAS -----

como os dados da estação são capturados? todas as variaveis, tempo de amostragem

Como eles são formatados para enviar ao ESP? tempo de coleta, métricas estatísticas, variáveis de interesse, vetor uart

estudar ainda essa formatação. dados de interesse: temp, umid, concentração de 5 compostos = 7 variáveis.

métricas de média, mínima e máxima, quantidade de amostras constante.

integração com o DMA

ESP recebe o pacote uart, formata pacote de envio wireless

recepcção de decodificação de pacote wireless na nuvem

integração com o servidor de aplicação

------ EM TERMOS GERAIS --------

A estação tem amostragem de dados a cada 4 segundos. Todos os dados da estação são amostrados e registrados localmente em cartão SD. A estação é focada em 5 sensores da fabricante Semeatech MODELO. Sensores de VOC, H2S, SO2, NH3 e CH3SH. Estes sensores eletroquímicos fornecem a saída digital da concentração do composto detectada e da temperatura e da umidade captadas pelo sensor que são utilizados nos algoritmos de compensação. 

O escopo do projeto prevê a correção e implementação das leituras via aprendizado de máquina com os dados brutos da estação na nuvem, e as variáveis esperadas pelo modelo são: concentração de cada composto e a média das leituras de temperatura e umidade de cada sensor.

Um protocolo de comunicação baseado em UART foi desenvolvido para tramitar dados entre o microcontrolador aquisitor e o microcontrolador de conectividade. CONFIGURAÇÔES DO PADRÂO UART. O protocolo segue o fluxo de chaves e respostas. Foi atribuído um código de um \textit{byte} para corresponder a cada um dos sensores. Assim, o ESP32 requisita os dados de um sensor ao STM32, o STM32 monta um pacote para transmissão serial e responded a solicitação do ESP32. Esse pacote de transmissão UART é padronizado no tamanho de 8 bytes, sendo o primeiro o ID do sensor, os seis subsequentes são as leituras do sensor, o oitavo com o valor da soma de verificação do pacote. No ESP32 é verificado se a reposta recebida era a esperada pela verificação do ID do sensor e a integridade é verificada a partir do último byte. Uma vez que o pacote serial é válido, há a acumulação dos dados na RAM do ESP32.

O STM32 acumula todas as leituras realizadas nas amostragens periódicas de 4 segundos até que o ESP32 solicite os pacotes. A cada ciclo de amostragem é acumulado o valor medido e verificado se valor máximo ou se valor mínimo. Um contador de amostras também é acrescido. Quando o ESP32 solicita as leituras do pacote. O STM32 divide os valos acumulados pelo contados de amostras para calcular a média de leituras obtidas entre as solicitações do ESP e separa os valores máximos e mínimos.

Os dados do pacote de comunicação serial é Id do sensor, MSB da média, LSB da média, MSB da máxima, LSB da máxima, MSB da mínima, LSB da mínima, (contador de amostras), checksum.

Com os dados das leituras no ESP32 é realizado a montagem do pacote para envio LoRaWAN.


No servidor de rede o pacote é decodificado através de um decodificado em JavaScript e os dados são salvos em um json. Depois conectado ao webhook explicar o IoT agente e por aí vai


Os dados são amostrados em painéis no grafana.
